\chapter{1988年上海卷}

\section{填空题}

\begin{problem}
方程 \(4^x-2^{x+1}-8=0\) 的解是\fillinblank{}.
\end{problem}

\begin{answer}
\(2\)
\end{answer}

\begin{problem}
\(y=\sqrt{x-1}\) 的反函数是\fillinblank{}.
\end{problem}

\begin{answer}
\(y=x^2+1\)（\(x\geqslant0\)）
\end{answer}

\begin{problem}
在极坐标系中，以 \(\left( \frac{a}{2},\frac{\pi}{2} \right)\) 为圆心，\(\frac{a}{2}\) 为半径的圆的方程为\fillinblank{}.
\end{problem}

\begin{answer}
\(\rho=a\sin\theta\)
\end{answer}

\begin{problem}
参数方程 \(\begin{cases}
x=1-2t,\\
y=2t^2-t
\end{cases}\)（\(t\) 为参数）化为普通方程是\fillinblank{}.
\end{problem}

\begin{answer}
\(y=\frac{x^2-x}{2}\)
\end{answer}

\begin{problem}
圆锥的侧面展开图是圆心角为 \(216^\circ\) 的扇形，则此圆锥的轴截面的半顶角的正弦是\fillinblank{}.
\end{problem}

\begin{answer}
\(\frac{3}{5}\)
\end{answer}

\begin{problem}
\(\displaystyle\lim_{n\to\infty}\frac{(-2)^{n+1}}{1-2+4-\cdots+(-2)^{n-1}}=\)\fillinblank{}.
\end{problem}

\begin{answer}
\(6\)
\end{answer}

\begin{problem}
函数 \(y=\cos\left( \frac{\pi}{2}x \right)\cos\left[ \frac{\pi}{2}(x-1) \right]\) 的最小正周期是\fillinblank{}.
\end{problem}

\begin{answer}
\(2\)
\end{answer}

\begin{problem}
在 \((1+x)^n\) 的展开式中，若第三项和第六项的系数相等，则 \(n=\)\fillinblank{}.
\end{problem}

\begin{answer}
\(7\)
\end{answer}

\begin{problem}
到椭圆 \(\frac{x^2}{25}+\frac{y^2}{9}=1\) 右焦点的距离与直线 \(x=6\) 的距离相等的点的轨迹方程是\fillinblank{}.
\end{problem}

\begin{answer}
\(y^2=-4(x-5)\)
\end{answer}

\begin{problem}
从 \(6\) 个运动员中选出 \(4\) 人参加 \(4\times100\text{ 米}\)接力赛，如果甲，乙两人都不能跑第一棒，那末共有\fillinblank{} 种不同的参赛方案（要求结果用数字表达）.
\end{problem}

\begin{answer}
\(240\)
\end{answer}

\section{单选题}

\begin{problem}
方程 \(x^4-y^4-4x^2+4y^2=0\) 所表示的曲线是（\quad）

\choices
    {双曲线和一个圆}
    {两条相交直线}
    {两条平行直线和一个圆}
    {两条相交直线和一个圆}
\end{problem}

\begin{answer}
D
\end{answer}

\begin{problem}
在 \(\left[ -1,\frac{3}{2} \right]\) 上与函数 \(y=x\) 相同的函数是（\quad）

\choices
    {\(y=\arccos(\cos x)\)}
    {\(y=\arcsin(\sin x)\)}
    {\(y=\sin(\arcsin x)\)}
    {\(y=\cos(\arccos x)\)}
\end{problem}

\begin{answer}
B
\end{answer}

\begin{problem}
函数 \(y=\frac{1+a^x}{1-a^x}\)（\(a>0\)，\(a\ne1\)）（\quad）

\choices
    {是奇函数}
    {是偶函数}
    {既是奇函数，又是偶函数}
    {是非奇非偶函数}
\end{problem}

\begin{answer}
A
\end{answer}

\begin{problem}
\(|x|\leqslant2\) 是 \(|x+1|<1\) 的（\quad）

\choices
    {必要不充分条件}
    {充分不必要条件}
    {充分必要条件}
    {既不充分也不必要条件}
\end{problem}

\begin{answer}
A
\end{answer}

\begin{problem}
下列命题中错误的是（\quad）

\choices
    {若一直线垂直于一平面，则此直线必垂直于这平面上所有直线}
    {若一个平面通过另一个平面的一条垂线，则这两个平面互相垂直}
    {若一直线垂直于一个平面的一条垂线，则此直线平行于这个平面}
    {若平面内的一条直线和这个平面的一条斜线的射影垂直，则它也和这条斜线垂直}
\end{problem}

\begin{answer}
C
\end{answer}

\begin{problem}
若 \(\sin\alpha=\frac{5}{13}\)，\(\alpha\) 为第二象限角，则 \(\tan\frac{\alpha}{2}\) 的值为（\quad）

\choices
    {\(\frac{1}{5}\)}
    {\(5\)}
    {\(-5\)}
    {\(-\frac{1}{5}\)}
\end{problem}

\begin{answer}
B
\end{answer}

\begin{problem}
在半径为 \(a\) 的球中，高为 \(\frac{2}{3}a\) 的球冠面积与整个球面面积之和为（\quad）

\choices
    {\(\frac{8}{3}\pi a^2\)}
    {\(\frac{10}{3}\pi a^2\)}
    {\(\frac{14}{3}\pi a^2\)}
    {\(\frac{16}{3}\pi a^2\)}
\end{problem}

\begin{answer}
D
\end{answer}

\begin{problem}
下列函数图象中正确的是（\quad）

\choices
    {\choicebitmap[width=0.15\paperwidth]{img/1988/shanghai/q18_optA.png}}
    {\choicebitmap[width=0.15\paperwidth]{img/1988/shanghai/q18_optB.png}}
    {\choicebitmap[width=0.15\paperwidth]{img/1988/shanghai/q18_optC.png}}
    {\choicebitmap[width=0.15\paperwidth]{img/1988/shanghai/q18_optD.png}}
\end{problem}

\begin{answer}
B
\end{answer}

\begin{problem}
下列关系中正确的是（\quad）

\choices
    {\(\left( \frac{1}{2} \right)^{\frac{2}{3}}<\left( \frac{1}{5} \right)^{\frac{2}{3}}<\left( \frac{1}{2} \right)^{\frac{1}{3}}\)}
    {\(\left( \frac{1}{2} \right)^{\frac{1}{3}}<\left( \frac{1}{2} \right)^{\frac{2}{3}}<\left( \frac{1}{5} \right)^{\frac{2}{3}}\)}
    {\(\left( \frac{1}{5} \right)^{\frac{1}{3}}<\left( \frac{1}{2} \right)^{\frac{1}{3}}<\left( \frac{1}{2} \right)^{\frac{2}{3}}\)}
    {\(\left( \frac{1}{5} \right)^{\frac{2}{3}}<\left( \frac{1}{2} \right)^{\frac{2}{3}}<\left( \frac{1}{2} \right)^{\frac{1}{3}}\)}
\end{problem}

\begin{answer}
D
\end{answer}

\begin{problem}
方程 \(\sin x=\lg x\) 的实数根有（\quad）

\choices
    {\(1\) 个}
    {\(2\) 个}
    {\(3\) 个}
    {无穷多个}
\end{problem}

\begin{answer}
C
\end{answer}

\section{解答题}

\begin{problem}
已知直线 \(l_1:3x-y+12=0\)，\(l_2:3x+2y-6=0\). 求 \(l_1\)，\(l_2\) 和 \(y\) 轴所围成的三角形面积.
\end{problem}

\begin{solution}
直线 \(l_1\)，\(l_2\) 在 \(y\) 轴上的截距分别为 \(12\) 和 \(3\)，故它们在 \(y\) 轴上所截得的线段长为 \(9\). 联立 \(3x-y+12=0\)，\(3x+2y-6=0\) 得 \(x=-2\). 所以所求三角形的面积为 \(S_{\triangle}=\frac{1}{2}\times9\times2=9\).
\end{solution}

\begin{problem}
设 \(A=\left\{ x\mid x\geqslant\frac{1}{x},x\in\R \right\}\)，\(B=\left\{ x\mid \sqrt{2x+1}<3,x\in\R \right\}\). 求 \(D=A\cap B\).
\end{problem}

\begin{solution}
由 \(x\geqslant\frac{1}{x}\) 得 \(A=\left\{ x\mid -1\leqslant x<0,\text{或 }x\geqslant1,\ x\in\R \right\}\). 由 \(\sqrt{2x+1}<3\) 得 \(B=\left\{ x\mid -\frac{1}{2}\leqslant x<4,\ x\in\R \right\}\). 因此
\[D=A\cap B=\left\{ x\mid -1\leqslant x<0,\text{或 }x\geqslant1 \right\}\cap\left\{ x\mid -\frac{1}{2}\leqslant x<4 \right\}=\left\{ x\mid -\frac{1}{2}\leqslant x<0,\text{或 }1\leqslant x<4 \right\}.\]
\end{solution}

\begin{problem}
一个圆锥容器和一个圆柱形容器，它们的轴截面尺寸如图. 两容器内盛有的液体的体积正好相等，且液面高度 \(h\) 正好相同. 求 \(h\).

\bitmapfigure[width=0.58\linewidth]{img/1988/shanghai/q23_fig1.png}
\end{problem}

\begin{solution}
圆锥形容器的轴截面的底角为 \(45^\circ\)，故液面半径 \(r=h\). 所以圆锥形容器内液体体积为 \(V_{\text{锥}}=\frac{1}{3}\pi r^2h=\frac{\pi}{3}h^3\). 圆柱形容器内液体体积为 \(V_{\text{柱}}=\pi\left( \frac{a}{2} \right)^2h=\frac{\pi a^2h}{4}\). 由题意 \(V_{\text{锥}}=V_{\text{柱}}\)，得 \(\frac{\pi}{3}h^3=\frac{\pi a^2h}{4}\)，解得 \(h=\frac{\sqrt{3}}{2}a\).
\end{solution}

\begin{problem}
设方程 \(x^2+3\sqrt{3}x+4=0\) 的两个实根为 \(x_1\)，\(x_2\). 记 \(\alpha=\arctan x_1\)，\(\beta=\arctan x_2\). 求 \(\alpha+\beta\).
\end{problem}

\begin{solution}
由韦达定理得 \(x_1+x_2=-3\sqrt{3}\)，\(x_1x_2=4\). 因为 \(\tan\alpha=x_1\)，\(\tan\beta=x_2\)，所以
\[
\tan(\alpha+\beta)=\frac{\tan\alpha+\tan\beta}{1-\tan\alpha\tan\beta}=\frac{x_1+x_2}{1-x_1x_2}=\sqrt{3}
\]
又因为 \(x_1x_2>0\) 且 \(x_1+x_2<0\)，所以 \(x_1<0\)，\(x_2<0\)，从而 \(\alpha\)，\(\beta\in\left( -\frac{\pi}{2},0 \right)\)，故 \(\alpha+\beta\in(-\pi,0)\). 因此 \(\alpha+\beta=-\frac{2\pi}{3}\).
\end{solution}

\begin{problem}
设 \(y=\log_{\frac{1}{2}}\left[ a^{2x}+2(ab)^x-b^{2x}+1 \right]\)（\(a>0\)，\(b>0\)）. 求使 \(y\) 为负值的 \(x\) 的取值范围.
\end{problem}

\begin{solution}
要使 \(y<0\)，只要 \(a^{2x}+2(ab)^x-b^{2x}>0\)，即 \(b^{2x}\left[ \left( \frac{a}{b} \right)^{2x}+2\left( \frac{a}{b} \right)^x-1 \right]>0\). 由因式分解得 \(b^{2x}\left[ \left( \frac{a}{b} \right)^x+1+\sqrt{2} \right]\left[ \left( \frac{a}{b} \right)^x+1-\sqrt{2} \right]>0\). 因为 \(b^{2x}>0\)，\(\left( \frac{a}{b} \right)^x+1+\sqrt{2}>0\)，故只需 \(\left( \frac{a}{b} \right)^x>\sqrt{2}-1\). 两端取对数得 \(x\lg\dfrac{a}{b}>\lg(\sqrt{2}-1)\).

当 \(a>b>0\) 时，\(\lg\frac{a}{b}>0\)，由上式得 \(x>\dfrac{\lg(\sqrt{2}-1)}{\lg\frac{a}{b}}=\log_{\frac{a}{b}}(\sqrt{2}-1)\). 当 \(b>a>0\) 时，\(\lg\frac{a}{b}<0\)，由上式得 \(x<\dfrac{\lg(\sqrt{2}-1)}{\lg\frac{a}{b}}=\log_{\frac{a}{b}}(\sqrt{2}-1)\). 当 \(a=b>0\) 时，\(\lg\frac{a}{b}=0\)，而 \(\lg(\sqrt{2}-1)<0\)，所以不等式恒成立，此时 \(x\) 为一切实数.
\end{solution}

\begin{problem}
在棱长都相等的四面体 \(A-BCD\) 中，\(E\)，\(F\) 分别为棱 \(AD\)，\(BC\) 的中点，连接 \(AF\)，\(CE\)，如图所示.

\begin{enumerate}
\item 求异面直线 \(AF\)，\(CE\) 所成角的大小；
\item 求 \(CE\) 与底面 \(BCD\) 所成角的大小（要求角的大小用反三角函数表示）.
\end{enumerate}

\bitmapfigure[width=0.34\linewidth]{img/1988/shanghai/q26_fig1.png}
\end{problem}

\begin{solution}
\begin{enumerate}
\item 连接 \(DF\)，设 \(G\) 为 \(DF\) 的中点，连接 \(GE\). 因为 \(AE=ED\)，\(FG=GD\)，所以 \(EG\parallel AF\). 故 \(\angle GEC\) 是异面直线 \(AF\) 与 \(CE\) 所成的角.
设棱长为 \(a\). 则 \(AF=CE=\dfrac{\sqrt{3}}{2}a\)，\(EG=\dfrac{1}{2}AF=\dfrac{\sqrt{3}}{4}a\).
连接 \(GC\)，则
\[
GC=\sqrt{FC^2+FG^2}=\sqrt{\frac{a^2}{4}+\frac{3a^2}{16}}=\frac{\sqrt{7}}{4}a
\]
由余弦定理得 \(\cos\angle GEC=\dfrac{EG^2+CE^2-GC^2}{2EG\cdot CE}=\dfrac{2}{3}\)，所以 \(\angle GEC=\arccos\dfrac{2}{3}\)
即异面直线 \(AF\) 与 \(CE\) 所成角的大小为 \(\arccos\frac{2}{3}\).

\item 作 \(EK\perp DF\)，垂足为 \(K\)，连接 \(KC\). 因为 \(AF\perp BC\)，\(DF\perp BC\)，所以 \(BC\perp\) 平面 \(ADF\). 又 \(BC\subset\) 平面 \(BCD\)，故平面 \(ADF\perp\) 平面 \(BCD\). 于是 \(EK\perp\) 平面 \(BCD\)，从而 \(\angle KCE\) 是 \(CE\) 与底面 \(BCD\) 所成的角.
在 \(\triangle AFD\) 中，连接 \(EF\)，且 \(EF\perp AD\). 则
\[
EF=\sqrt{FD^2-ED^2}=\sqrt{\frac{3a^2}{4}-\frac{a^2}{4}}=\frac{\sqrt{2}}{2}a
\]
又因为 \(EK\cdot FD=EF\cdot ED\)，所以
\(EK=\frac{EF\cdot ED}{FD}=\frac{\frac{\sqrt{2}}{2}a\cdot\frac{a}{2}}{\frac{\sqrt{3}}{2}a}=\frac{\sqrt{6}}{6}a\).
于是
\[
\sin\angle KCE=\frac{EK}{CE}=\frac{\frac{\sqrt{6}}{6}a}{\frac{\sqrt{3}}{2}a}=\frac{\sqrt{2}}{3}
\]
所以 \(\angle KCE=\arcsin\dfrac{\sqrt{2}}{3}\)
即 \(CE\) 与底面 \(BCD\) 所成角的大小为 \(\arcsin\frac{\sqrt{2}}{3}\).
\end{enumerate}
\end{solution}

\begin{problem}
复数 \(z_1\)，\(z_2\)，\(z_3\) 的辐角分别为 \(\alpha\)，\(\beta\)，\(\gamma\)，又 \(|z_1|=1\)，\(|z_2|=k\)，\(|z_3|=2-k\)，且 \(z_1+z_2+z_3=0\). 问 \(k\) 各为何值时，\(\cos(\beta-\gamma)\) 分别取得最大值和最小值，并求出最大值和最小值.
\end{problem}

\begin{solution}
设 \(z_1=\cos\alpha+\i\sin\alpha\)，\(z_2=k\cos\beta+\i k\sin\beta\)，\(z_3=(2-k)\cos\gamma+\i(2-k)\sin\gamma\). 由 \(z_1+z_2+z_3=0\) 得
\(\cos\alpha+k\cos\beta+(2-k)\cos\gamma=0\)，\(\sin\alpha+k\sin\beta+(2-k)\sin\gamma=0\)，
即 \(\cos\alpha=(k-2)\cos\gamma-k\cos\beta\)，\(\sin\alpha=(k-2)\sin\gamma-k\sin\beta\).

将上面两式两端平方后相加，得 \(1=(k-2)^2+k^2-2(k-2)k\cos(\gamma-\beta)\). 由题设可知 \(k\ne0\)，\(k\ne2\)，所以 \(\cos(\beta-\gamma)=\dfrac{(k-2)^2+k^2-1}{2k(k-2)}=1+\dfrac{3}{2(k-1)^2-2}\).

又因为 \(\left| \cos(\beta-\gamma) \right|\leqslant1\)，故 \(-2\leqslant\dfrac{3}{2(k-1)^2-2}\leqslant0\)，解得 \(\dfrac{1}{2}\leqslant k\leqslant\dfrac{3}{2}\).

由上式可知，当 \(k=1\) 时，\(\left[ \cos(\beta-\gamma) \right]_{\max}=-\dfrac{1}{2}\)；当 \(k=\dfrac{1}{2}\) 或 \(k=\dfrac{3}{2}\) 时，\(\left[ \cos(\beta-\gamma) \right]_{\min}=-1\). 所以 \(\cos(\beta-\gamma)\) 的最大值为 \(-\dfrac{1}{2}\)，在 \(k=1\) 时取得；最小值为 \(-1\)，在 \(k=\dfrac{1}{2}\) 或 \(k=\dfrac{3}{2}\) 时取得.
\end{solution}

\begin{problem}
在双曲线 \(\frac{y^2}{12}-\frac{x^2}{13}=1\) 的一支上不同的三点 \(A(x_1,y_1)\)，\(B(\sqrt{26},6)\)，\(C(x_2,y_2)\) 与焦点 \(F(0,5)\) 的距离成等差数列.

\begin{enumerate}
\item 求 \(y_1+y_2\)；
\item 证明线段 \(AC\) 的垂直平分线经过某一定点，并求该定点的坐标.
\end{enumerate}
\end{problem}

\begin{solution}
\begin{enumerate}
\item 法一：双曲线的实半轴 \(a=2\sqrt{3}\)，虚半轴 \(b=\sqrt{13}\)，半焦距 \(c=5\). 由双曲线的定义可得 \(\left| AF \right|=\dfrac{c}{a}y_1-a=\dfrac{5\sqrt{3}}{6}y_1-2\sqrt{3}\)，同理 \(\left| CF \right|=\dfrac{5\sqrt{3}}{6}y_2-2\sqrt{3}\). 又 \(\left| BF \right|=\sqrt{(\sqrt{26}-0)^2+(6-5)^2}=3\sqrt{3}\). 因为 \(\left| AF \right|\)，\(\left| BF \right|\)，\(\left| CF \right|\) 成等差数列，所以 \(\left| AF \right|+\left| CF \right|=2\left| BF \right|\)，于是 \(\left( \dfrac{5\sqrt{3}}{6}y_1-2\sqrt{3} \right)+\left( \dfrac{5\sqrt{3}}{6}y_2-2\sqrt{3} \right)=6\sqrt{3}\)，解得 \(y_1+y_2=12\).

法二：因为 \(\left| BF \right|=\sqrt{(\sqrt{26}-0)^2+(6-5)^2}=\sqrt{26+1}=3\sqrt{3}\). 又 \(\left| AF \right|=\sqrt{x_1^2+(y_1-5)^2}\). 而 \(A(x_1,y_1)\) 在双曲线上，所以 \(x_1^2=\dfrac{13y_1^2-156}{12}\). 从而
\[
\left| AF \right|^2=\frac{13y_1^2-156}{12}+y_1^2-10y_1+25=\frac{1}{12}(25y_1^2-120y_1+144)=\frac{1}{12}(5y_1-12)^2
\]
由题设可知 \(y_1\geqslant\sqrt{12}\)，所以 \(5y_1-12>0\). 因此 \(\left| AF \right|=\dfrac{1}{\sqrt{12}}(5y_1-12)\). 同理可得 \(\left| CF \right|=\dfrac{1}{\sqrt{12}}(5y_2-12)\).
因为 \(\left| AF \right|+\left| CF \right|=2\left| BF \right|\)，所以 \(\dfrac{1}{\sqrt{12}}(5y_1-12+5y_2-12)=6\sqrt{3}\)，即得 \(y_1+y_2=12\).

\item 法一：线段 \(AC\) 的中点为 \(\left( \dfrac{x_1+x_2}{2},6 \right)\)，故它的垂直平分线 \(l\) 的方程为 \(y-6=-\dfrac{x_1-x_2}{y_1-y_2}\left( x-\dfrac{x_1+x_2}{2} \right)\). 由双曲线的一支关于 \(y\) 轴对称可知，若 \(l\) 经过一定点，则该定点必在 \(y\) 轴上. 设 \(l\) 与 \(y\) 轴的交点为 \(D(0,y_0)\)，则 \(y_0-6=\dfrac{x_1^2-x_2^2}{2(y_1-y_2)}\).
因为 \(A(x_1,y_1)\)，\(C(x_2,y_2)\) 在双曲线上，所以 \(x_1^2=\dfrac{13}{12}(y_1^2-12)\)，\(x_2^2=\dfrac{13}{12}(y_2^2-12)\). 因此
\[x_1^2-x_2^2=\frac{13}{12}(y_1^2-y_2^2)=\frac{13}{12}(y_1+y_2)(y_1-y_2)=13(y_1-y_2)\]
代入上式得 \(y_0-6=\dfrac{13}{2}\)，所以 \(y_0=\dfrac{25}{2}\).
因为 \(y_0\) 的值与 \(A\)，\(C\) 点的坐标无关，所以 \(D\) 为定点，即线段 \(AC\) 的垂直平分线经过定点 \(D\left( 0,\dfrac{25}{2} \right)\).

法二：线段 \(AC\) 的中点为 \(\left( \dfrac{x_1+x_2}{2},6 \right)\)，由题设可知 \(y_1\ne y_2\)，故它的垂直平分线 \(l\) 的方程为 \(y-6=-\dfrac{x_1-x_2}{y_1-y_2}\left( x-\dfrac{x_1+x_2}{2} \right)\). 由双曲线的一支关于 \(y\) 轴的对称性可知，若线段 \(AC\) 的垂直平分线经过某一定点，则该点必在 \(y\) 轴上，因此只需证 \(AC\) 的垂直平分线与 \(y\) 轴的交点为定点即可.
设 \(AC\) 的垂直平分线与 \(y\) 轴的交点为 \(D(0,y_0)\). 因为 \(\left| AD \right|=\left| CD \right|\)，所以 \(x_1^2+(y_1-y_0)^2=x_2^2+(y_2-y_0)^2\)，即
\[
\frac{13y_1^2-156}{12}+y_1^2-2y_1y_0+y_0^2=\frac{13y_2^2-156}{12}+y_2^2-2y_2y_0+y_0^2
\]
化简后得 \(25(y_1+y_2)(y_1-y_2)+24(y_2-y_1)y_0=0\). 由题设知 \(y_1\ne y_2\)，所以 \(25(y_1+y_2)-24y_0=0\)，解得 \(y_0=\dfrac{25\times12}{24}=\dfrac{25}{2}\). 因为 \(y_0\) 与 \(A\)，\(C\) 点的坐标无关，故 \(D\) 为定点，即线段 \(AC\) 的垂直平分线通过定点 \(D\left( 0,\dfrac{25}{2} \right)\).
\end{enumerate}
\end{solution}
